\section{Claim-Relevant Sensitivities and Verification Tables}
\label{app:sensitivities}

This appendix reports the verification detail needed to judge whether the main conclusions depend on an analysis choice or an unstable aggregate.
Section~\ref{sec:results} owns the decisive observations and scientific verdicts; Appendix~\ref{app:evidence-provenance} owns the visible evidence-record crosswalk and the verified subset of path--hash bindings.
Every value below traces to its owning evidence record, including values whose artifacts are not listed in that verified subset.

The tables preserve the distinction between inference units.
Participant-level effects support inference across the sampled biological units under the assumptions in Appendix~\ref{app:inference}.
Held-out blocks diagnose dependence on the fixed stimulus partition, and optimization seeds diagnose the stability of a fixed training procedure.
Neither blocks nor seeds create additional biological or model-population replicates.

\subsection{Claim-relevant analysis choices}

An analysis choice can affect a thesis claim in four ways.
An \emph{interpretation-changing} analysis replaces an earlier estimand or inference unit.
A \emph{scope-bounding} sensitivity narrows what the primary result supports without reversing its frozen decision rule.
A \emph{confirmatory} check leaves the scoped interpretation unchanged.
A \emph{limitation or correction} removes authority from a design element rather than generating a more favorable result.
Tables~\ref{tab:claim-changing-sensitivities} and~\ref{tab:confirmatory-sensitivities} apply these labels to the analyses that matter for the thesis claims.

\begingroup
\footnotesize
\setlength{\tabcolsep}{3pt}
\renewcommand{\arraystretch}{1.03}
\begin{longtable}{@{}p{28mm}p{31mm}p{22mm}p{51mm}@{}}
  \caption{Interpretation-changing and scope-bounding analyses.}
  \label{tab:claim-changing-sensitivities}\\
  \toprule
  Analysis choice & Scientific role & Status & Effect on interpretation \\
  \midrule
  \endfirsthead
  \multicolumn{4}{c}{\tablename~\thetable{} (continued)} \\
  \toprule
  Analysis choice & Scientific role & Status & Effect on interpretation \\
  \midrule
  \endhead
  \bottomrule
  \endlastfoot

  Representation rank
& Distillation-headroom analysis
  & Scope-bounding
  & The ordering of the model arms survived, but absolute magnitudes and some signs changed. Headroom is conditional on the declared rank.\evd{E003} \\

  \specialrule{0.15pt}{0.3ex}{0.3ex}

  Auxiliary weight and language quality
  & Participant-averaged recorded-response intervention
  & Scope-bounding
  & Near-quality-matched settings remained uncertain across folds. The clearest fold-positive setting incurred a substantial language-quality cost.\evd{E005} \\

  \specialrule{0.15pt}{0.3ex}{0.3ex}

  Quality band and model family
& Distillation-headroom analysis
  & Scope-bounding
  & The full-range association was \result{e015_full_corr}; within the capable-model band it was \result{e015_capable_corr}, with interval \result{e015_capable_ci}. Language quality remains a required control, not a deterministic explanation.\evd{E015} \\

  \specialrule{0.15pt}{0.3ex}{0.3ex}

  Response construction and inference unit
  & Participant-specific recorded-response intervention
  & Interpretation-changing
  & Participant-specific targets and participant-first inference answer a different question from an averaged target and fold-level summary. The change removed support for a participant-general gain.\evd{E008} \\

  \specialrule{0.15pt}{0.3ex}{0.3ex}

  Participant and subgroup deletion
& Saved-student biological-transfer assay and 50-component linear target-projection assay
  & Scope-bounding
& Removing the predeclared highest-baseline participant reversed the relative transfer point estimate. Every leave-one-participant-out linear target-projection mean remained negative.\evd{E025}\evd{E030} \\

  \specialrule{0.15pt}{0.3ex}{0.3ex}

  Held-out stimulus block
& Controlled naturalistic voxelwise predictivity assay; recorded-response intervention; saved-student biological-transfer assay; 50-component linear target-projection assay
  & Scope-bounding
& The naturalistic assay remains a one-participant result. No leave-one-block-out intervention analysis produced a positive participant-mean contrast that changed the participant-level disposition, and the linear target-projection result held under every block deletion.\evd{E006}\evd{E008}\evd{E025}\evd{E030} \\

  \specialrule{0.15pt}{0.3ex}{0.3ex}

  Controlled-predictivity contrast by seed
  & Practical-utility evaluation
  & Scope-bounding
  & Only \result{e009_positive_seed_signs} technical-seed signs were positive. The utility analysis therefore follows a seed-unstable prerequisite rather than a reliably induced target-specific change.\evd{E009} \\
\end{longtable}
\endgroup

\begin{table}[H]
  \centering
  \footnotesize
  \setlength{\tabcolsep}{4pt}
  \renewcommand{\arraystretch}{1.03}
  \caption{Confirmatory checks and design limitations.}
  \label{tab:confirmatory-sensitivities}
  \begin{tabularx}{\textwidth}{@{}>{\raggedright\arraybackslash}p{34mm}>{\raggedright\arraybackslash}p{39mm}>{\raggedright\arraybackslash}p{25mm}>{\raggedright\arraybackslash}X@{}}
    \toprule
    Analysis choice & Scientific role & Status & Effect on interpretation \\
    \midrule

    Layer, nuisance set, and target aggregation
& Saved-student biological-transfer assay and 50-component linear target-projection assay
    & Confirmatory
    & Prespecified nearby choices did not overturn the primary rules. The 50-component linear target projection did not meet its continuation rule; retention and composed-path quantities remain diagnostics of that pathway.\evd{E025}\evd{E030} \\

    \specialrule{0.15pt}{0.3ex}{0.3ex}

    Saved-student optimization seed
    & Synthetic-target recovery assay; retained-student movement audit; saved-student target-retention analysis; saved-student biological-transfer assay; composed-path diagnostic
    & Confirmatory
    & Seed agreement can establish reproducible optimization or manipulation, but it cannot substitute for participant agreement.\evd{E016}\evd{E025}\evd{E030} \\

    \specialrule{0.15pt}{0.3ex}{0.3ex}

    Saved controls audited after training
    & Comparator-adequacy audit and attribution assessment
    & Limitation or correction
    & The saved block-permutation sensitivity is defective, and the text-derived auxiliary control fails pre-outcome and KD-baseline comparability. Neither can identify an effect of brain-response content.\evd{E016}\evd{E026} \\

    \bottomrule
  \end{tabularx}
\end{table}

\begin{table}[H]
  \centering
  \small
  \caption{Corrected classification of the frozen comparator audit.}
  \label{tab:e026-corrected-audit}
  \begin{tabularx}{\textwidth}{@{}l r r r X@{}}
    \toprule
    Role & Total & Pass or within & Fail or outside & Interpretation \\
    \midrule
    Identification checks & 37 & 20 & 16 & One additional check is unresolved; failures include KD-baseline recovery and headroom, covariance geometry, effective rank, and nuisance predictability. \\
    Post-training diagnostics & 18 & 4 & 14 & Parameter, Frobenius, and \ac{cka} margins describe realized intervention differences and do not determine comparator validity. \\
    \bottomrule
  \end{tabularx}
\end{table}

The original combined engineering count, \result{e026_engineering_split}, is retained only as provenance for the frozen audit.\evd{E026}

The interpretation-changing analysis altered the question being answered; it was not a more favorable reanalysis of the earlier averaged-target result.
The remaining sensitivities either narrowed the scope, confirmed stability near a frozen specification, or exposed a limitation.
None licenses replacing a prospective primary estimator or decision rule with a more favorable variant.

\subsection{Participant-level intervention contrasts}

The participant-specific recorded-response estimator first forms a paired target-specific contrast for participant \(u\), held-out block \(k\), and optimization seed \(s\), where \(P=5\) is the number of matched permuted-target draws.\evd{E008}
\begin{equation}
  d_{uks}
  = R^{2,\mathrm{unique}}_{\mathrm{recorded},uks}
  - \frac{1}{P}\sum_{p=1}^{P}R^{2,\mathrm{unique}}_{\mathrm{perm}(p),uks},
  \qquad
  e_u=\operatorname{median}_{k,s}d_{uks},
  \qquad
  \widehat{\theta}_{\mathrm{part}}=\frac{1}{U}\sum_{u=1}^{U}e_u.
  \label{eq:appendix-recorded-participant}
\end{equation}
Equation~\ref{eq:appendix-recorded-participant} subtracts the mean matched permuted-target score inside every seed--block cell, then aggregates the crossed cells within participant before averaging participants.
The resulting \(U=\result{e008_participant_count}\) participant effects, each based on \result{e008_cells_per_participant} paired cells, are the biological inputs to the cohort estimate.

\begin{table}[H]
  \centering
  \small
  \setlength{\tabcolsep}{12pt}
  \renewcommand{\arraystretch}{1.02}
  \caption{Participant-specific recorded-response intervention contrasts.}
  \label{tab:recorded-participant-values}
  \begin{tabular}{rc}
    \toprule
    \ac{uid} & Participant effect \(e_u\) \\
    \midrule
    797 & \result{e008_uid797} \\
    837 & \result{e008_uid837} \\
    841 & \result{e008_uid841} \\
    848 & \result{e008_uid848} \\
    856 & \result{e008_uid856} \\
    865 & \result{e008_uid865} \\
    875 & \result{e008_uid875} \\
    876 & \result{e008_uid876} \\
    880 & \result{e008_uid880} \\
    \midrule
    Mean & \result{e008_mean} \\
    \bottomrule
  \end{tabular}
\end{table}

The signs are mixed around a small cohort mean.
The table supports participant-first aggregation and heterogeneity inspection; it does not imply that every participant has exactly zero effect. \evd{E008}

The saved-student biological-transfer assay uses a different target and within-participant aggregator.
Let \(T_u\), \(X_u\), and \(K_u\) denote the participant-level scores for the synthetic-response target, text-derived auxiliary control, and ordinary \ac{kd}.
The three displayed contrasts satisfy
\begin{equation}
  \Delta^{T-X}_u=T_u-X_u,
  \quad
  \Delta^{T-K}_u=T_u-K_u,
  \quad
  \Delta^{X-K}_u=X_u-K_u,
  \quad
  \Delta^{T-X}_u=\Delta^{T-K}_u-\Delta^{X-K}_u.
  \label{eq:appendix-transfer-identity}
\end{equation}
Equation~\ref{eq:appendix-transfer-identity} shows why a positive relative contrast against the auxiliary control need not imply improvement over ordinary \ac{kd}: the auxiliary control may instead have moved downward.
Seeds and outer folds are averaged within each participant before these contrasts are formed.

\begin{table}[H]
  \centering
  \small
  \setlength{\tabcolsep}{8pt}
  \renewcommand{\arraystretch}{1.02}
  \caption{Participant-level saved-student biological-transfer contrasts.}
  \label{tab:transfer-participant-values}
  \begin{tabular}{rccc}
    \toprule
    \ac{uid} & \(\Delta^{T-X}_u\) & \(\Delta^{T-K}_u\) & \(\Delta^{X-K}_u\) \\
    \midrule
    797 & \result{e025_uid797_brain_text} & \result{e025_uid797_brain_kd} & \result{e025_uid797_text_kd} \\
    837 & \result{e025_uid837_brain_text} & \result{e025_uid837_brain_kd} & \result{e025_uid837_text_kd} \\
    841 & \result{e025_uid841_brain_text} & \result{e025_uid841_brain_kd} & \result{e025_uid841_text_kd} \\
    848 & \result{e025_uid848_brain_text} & \result{e025_uid848_brain_kd} & \result{e025_uid848_text_kd} \\
    856 & \result{e025_uid856_brain_text} & \result{e025_uid856_brain_kd} & \result{e025_uid856_text_kd} \\
    865 & \result{e025_uid865_brain_text} & \result{e025_uid865_brain_kd} & \result{e025_uid865_text_kd} \\
    875 & \result{e025_uid875_brain_text} & \result{e025_uid875_brain_kd} & \result{e025_uid875_text_kd} \\
    876 & \result{e025_uid876_brain_text} & \result{e025_uid876_brain_kd} & \result{e025_uid876_text_kd} \\
    880 & \result{e025_uid880_brain_text} & \result{e025_uid880_brain_kd} & \result{e025_uid880_text_kd} \\
    \midrule
    Mean & \result{e025_participant_delta} & \result{e025_tribe_kd} & \result{e025_textfeat_kd} \\
    \bottomrule
  \end{tabular}
\end{table}

The transfer contrasts are also heterogeneous.
The relative mean \(\Delta^{T-X}\) is small and positive, while the direct mean \(\Delta^{T-K}\) is \result{e025_tribe_kd}; removing the predeclared highest-baseline participant changes the relative point estimate to \result{e025_drop_highest}.
These rows explain why the relative comparator cannot replace the direct contrast or participant-first inference. \evd{E025}\evd{E026}

The two participant tables must not be pooled or compared row by row.
They use different targets, contrasts, and within-participant aggregators; the shared identifiers only show that they were evaluated in the same participant cohort.

\subsection{Participant-level mechanism diagnostics}

The target-projection diagnostic evaluates the declared 50-component linear pathway on the same sentence substrate; it is not a prerequisite for the independent direct-transfer contrast.
For each participant, \(A_u\) measures the target's unique response predictivity above nuisance, and \(Q_u\) compares that score with the frozen row-twin control.
These quantities compare target-above-nuisance predictivity with sentence-pairing specificity before the later target-retention analysis and composed-path diagnostic are interpreted.

\begin{table}[H]
  \centering
  \small
  \setlength{\tabcolsep}{10pt}
  \renewcommand{\arraystretch}{1.02}
  \caption{Participant-level linear target-projection diagnostics.}
  \label{tab:exact-target-participant-values}
  \begin{tabular}{rcc}
    \toprule
    \ac{uid} & Target above nuisance \(A_u\) & Aligned minus frozen twin \(Q_u\) \\
    \midrule
    797 & \result{e030_uid797_target_unique} & \result{e030_uid797_target_twin} \\
    837 & \result{e030_uid837_target_unique} & \result{e030_uid837_target_twin} \\
    841 & \result{e030_uid841_target_unique} & \result{e030_uid841_target_twin} \\
    848 & \result{e030_uid848_target_unique} & \result{e030_uid848_target_twin} \\
    856 & \result{e030_uid856_target_unique} & \result{e030_uid856_target_twin} \\
    865 & \result{e030_uid865_target_unique} & \result{e030_uid865_target_twin} \\
    875 & \result{e030_uid875_target_unique} & \result{e030_uid875_target_twin} \\
    876 & \result{e030_uid876_target_unique} & \result{e030_uid876_target_twin} \\
    880 & \result{e030_uid880_target_unique} & \result{e030_uid880_target_twin} \\
    \midrule
    Mean & \result{e030_target_unique} & \result{e030_target_twin} \\
    \bottomrule
  \end{tabular}
\end{table}

The mean target-above-nuisance score is \result{e030_target_unique}, and every leave-one-participant-out mean remains negative.
This first prerequisite misses its frozen practical rule.
The aligned-minus-frozen-twin values have substantial effects of both signs, so that paired specificity question remains unresolved rather than contradicting the first result. \evd{E030}

The later diagnostics ask whether the part of the target that a student predicts overlaps with the target direction that predicts participant responses.
For participant \(u\), \(O_u\) is the absolute composed student-to-target-to-response overlap, \(D_u\) is its brain-target-minus-\ac{kd} increment, and \(J_u\) is the aligned-minus-frozen-twin increment.
All three average optimization seeds and outer folds within participant before participant-level inspection.

\begin{table}[H]
  \centering
  \small
  \setlength{\tabcolsep}{8pt}
  \renewcommand{\arraystretch}{1.02}
  \caption{Participant-level composed-path quantities.}
  \label{tab:composed-participant-values}
  \begin{tabular}{rccc}
    \toprule
    \ac{uid} & Absolute overlap \(O_u\) & Synthetic response minus \ac{kd} \(D_u\) & Aligned minus twin \(J_u\) \\
    \midrule
    797 & \result{e030_uid797_overlap_abs} & \result{e030_uid797_overlap_kd} & \result{e030_uid797_overlap_twin} \\
    837 & \result{e030_uid837_overlap_abs} & \result{e030_uid837_overlap_kd} & \result{e030_uid837_overlap_twin} \\
    841 & \result{e030_uid841_overlap_abs} & \result{e030_uid841_overlap_kd} & \result{e030_uid841_overlap_twin} \\
    848 & \result{e030_uid848_overlap_abs} & \result{e030_uid848_overlap_kd} & \result{e030_uid848_overlap_twin} \\
    856 & \result{e030_uid856_overlap_abs} & \result{e030_uid856_overlap_kd} & \result{e030_uid856_overlap_twin} \\
    865 & \result{e030_uid865_overlap_abs} & \result{e030_uid865_overlap_kd} & \result{e030_uid865_overlap_twin} \\
    875 & \result{e030_uid875_overlap_abs} & \result{e030_uid875_overlap_kd} & \result{e030_uid875_overlap_twin} \\
    876 & \result{e030_uid876_overlap_abs} & \result{e030_uid876_overlap_kd} & \result{e030_uid876_overlap_twin} \\
    880 & \result{e030_uid880_overlap_abs} & \result{e030_uid880_overlap_kd} & \result{e030_uid880_overlap_twin} \\
    \midrule
    Mean & \result{e030_overlap_abs} & \result{e030_overlap_delta} & \result{e030_overlap_twin} \\
    \bottomrule
  \end{tabular}
\end{table}

Absolute composed overlap remains unresolved.
The two increments are smaller and have mixed signs.
Within the declared linear pathway, these downstream diagnostics cannot repair the target-projection result or identify its cause. They do not constrain the independent direct-transfer contrast. \evd{E030}

\subsection{Technical-seed diagnostics}

Optimization seeds test whether a fixed training procedure behaves reproducibly.
They do not support biological generalization, because the same participants and stimuli are reused within each seed.

The brain-specific interpretation of the practical-utility evaluation depends on a reliable prerequisite contrast in controlled brain predictivity between the recorded-response arm and its matched permuted-response control.
Table~\ref{tab:utility-seed-values} reports the complete seed vector used to judge that prerequisite.

\begin{table}[H]
  \centering
  \small
  \setlength{\tabcolsep}{12pt}
  \renewcommand{\arraystretch}{1.02}
  \caption{Technical-seed values for the practical-utility prerequisite.}
  \label{tab:utility-seed-values}
  \begin{tabular}{rc}
    \toprule
    Seed & Target minus permuted-target representation contrast \\
    \midrule
    0 & \result{e009_seed0} \\
    1 & \result{e009_seed1} \\
    2 & \result{e009_seed2} \\
    3 & \result{e009_seed3} \\
    4 & \result{e009_seed4} \\
    5 & \result{e009_seed5} \\
    6 & \result{e009_seed6} \\
    7 & \result{e009_seed7} \\
    \midrule
    Mean & \result{e009_fulcrum_mean} \\
    Median & \result{e009_fulcrum_median} \\
    \bottomrule
  \end{tabular}
\end{table}

The signs are split, and the mean is strongly influenced by seed~0.
The seed rows are rounded independently to three decimal places; the displayed mean and median are computed from artifact-precision values.
The complete vector documents mixed signs and seed-0 sensitivity; it does not establish a null or equivalence. \evd{E009}

For the saved-student biological-transfer assay, each seed row averages participants and outer folds.
The three columns use the same contrasts as Equation~\ref{eq:appendix-transfer-identity}.

\begin{table}[H]
  \centering
  \small
  \setlength{\tabcolsep}{8pt}
  \renewcommand{\arraystretch}{1.02}
  \caption{Technical-seed values for saved-student biological transfer.}
  \label{tab:transfer-seed-values}
  \begin{tabular}{rccc}
    \toprule
    Seed & \(\Delta^{T-X}\) & \(\Delta^{T-K}\) & \(\Delta^{X-K}\) \\
    \midrule
    0 & \result{e025_seed0_brain_text} & \result{e025_seed0_brain_kd} & \result{e025_seed0_text_kd} \\
    1 & \result{e025_seed1_brain_text} & \result{e025_seed1_brain_kd} & \result{e025_seed1_text_kd} \\
    2 & \result{e025_seed2_brain_text} & \result{e025_seed2_brain_kd} & \result{e025_seed2_text_kd} \\
    3 & \result{e025_seed3_brain_text} & \result{e025_seed3_brain_kd} & \result{e025_seed3_text_kd} \\
    4 & \result{e025_seed4_brain_text} & \result{e025_seed4_brain_kd} & \result{e025_seed4_text_kd} \\
    5 & \result{e025_seed5_brain_text} & \result{e025_seed5_brain_kd} & \result{e025_seed5_text_kd} \\
    \midrule
    Mean & \result{e025_participant_delta} & \result{e025_tribe_kd} & \result{e025_textfeat_kd} \\
    \bottomrule
  \end{tabular}
\end{table}

Before participant responses were inspected, the synthetic-target recovery, retained-student movement, and language-quality criteria were satisfied.
The largest saved-arm language-quality difference, \result{e025_quality_max_gap}, remained within the frozen \result{e025_quality_margin} margin.
The seed table therefore checks an interpretable manipulation while leaving biological inference to the participant table. \evd{E016}\evd{E025}

The exact-target seed diagnostics separate absolute target extraction from improvement over seed-matched ordinary \ac{kd} and from composed-path increments.
Within each seed, target extraction and incremental retention aggregate held-out target recovery; the two overlap increments average participants and outer folds.

\begin{table}[H]
  \centering
  \small
  \setlength{\tabcolsep}{10pt}
  \renewcommand{\arraystretch}{1.02}
  \caption{Technical-seed target-extraction and retention values.}
  \label{tab:exact-target-seed-values}
  \begin{tabular}{rcc}
    \toprule
    Seed & Target extraction \(M_s\) & Incremental retention \(B_s\) \\
    \midrule
    0 & \result{e030_seed0_target_extract} & \result{e030_seed0_retention_delta} \\
    1 & \result{e030_seed1_target_extract} & \result{e030_seed1_retention_delta} \\
    2 & \result{e030_seed2_target_extract} & \result{e030_seed2_retention_delta} \\
    3 & \result{e030_seed3_target_extract} & \result{e030_seed3_retention_delta} \\
    4 & \result{e030_seed4_target_extract} & \result{e030_seed4_retention_delta} \\
    5 & \result{e030_seed5_target_extract} & \result{e030_seed5_retention_delta} \\
    \midrule
    Mean & \result{e030_target_extract} & \result{e030_retention_delta} \\
    \bottomrule
  \end{tabular}
\end{table}

\begin{table}[H]
  \centering
  \small
  \setlength{\tabcolsep}{10pt}
  \renewcommand{\arraystretch}{1.02}
  \caption{Technical-seed composed-path increments.}
  \label{tab:composed-seed-values}
  \begin{tabular}{rcc}
    \toprule
    Seed & Synthetic response minus \ac{kd} \(D_s\) & Aligned minus twin \(J_s\) \\
    \midrule
    0 & \result{e030_seed0_overlap_kd} & \result{e030_seed0_overlap_twin} \\
    1 & \result{e030_seed1_overlap_kd} & \result{e030_seed1_overlap_twin} \\
    2 & \result{e030_seed2_overlap_kd} & \result{e030_seed2_overlap_twin} \\
    3 & \result{e030_seed3_overlap_kd} & \result{e030_seed3_overlap_twin} \\
    4 & \result{e030_seed4_overlap_kd} & \result{e030_seed4_overlap_twin} \\
    5 & \result{e030_seed5_overlap_kd} & \result{e030_seed5_overlap_twin} \\
    \midrule
    Mean & \result{e030_overlap_delta} & \result{e030_overlap_twin} \\
    \bottomrule
  \end{tabular}
\end{table}

Absolute target extraction is positive in every seed.
Incremental retention has mixed signs and an interval containing zero.
Both composed-path increments show the same uncertainty.
The separation shows that fitting the target in absolute terms does not establish improvement over ordinary \ac{kd} or transfer along the composed path. \evd{E030}

\begin{table}[H]
  \centering
  \small
  \caption{Primary retained-student movement values by technical seed.}
  \label{tab:e025-movement-seeds}
  \begin{tabular}{@{}r r r@{}}
    \toprule
    Seed & Centered relative-Frobenius & Linear-\ac{cka} distance \\
    \midrule
    0 & 0.186369 & 0.002633 \evd{E025} \\
    1 & 0.216294 & 0.005457 \evd{E025} \\
    2 & 0.298694 & 0.011778 \evd{E025} \\
    3 & 0.195493 & 0.003585 \evd{E025} \\
    4 & 0.197089 & 0.003925 \evd{E025} \\
    5 & 0.196519 & 0.003423 \evd{E025} \\
    \bottomrule
  \end{tabular}
\end{table}

The values use layer 6 and the same 1,999 untouched WikiText sentences in every arm; Equation~\ref{eq:movement-distances} defines both distances.\evd{E025}

No seed-only table is reported for the participant-specific recorded-response intervention.
Its primary participant statistic is a median over crossed seed--block cells, so a seed-only vector cannot reconstruct the nonlinear participant aggregate in Equation~\ref{eq:appendix-recorded-participant}.
Seed variability remains inside each participant value, while the shared-block summary diagnoses the other crossed axis. \evd{E008}

\begin{longtable}{@{}p{42mm}p{17mm}r r r r@{}}
\caption{Models in the cross-family language-quality analysis.}\label{tab:e015-model-panel}\\
\toprule
Model & Family & Parameters & Mid layer & Bits/byte & Unique \(R^2\) \\
\midrule
\endfirsthead
\toprule
Model & Family & Parameters & Mid layer & Bits/byte & Unique \(R^2\) \\
\midrule
\endhead
Mistral-7B-v0.1 & Mistral & 7.24B & 16 & 1.135 & +0.0438 \evd{E015} \\
Qwen2.5-7B & Qwen & 7.62B & 14 & 1.152 & +0.0235 \evd{E015} \\
Llama-3.2-3B & Llama & 3.21B & 14 & 1.163 & +0.0464 \evd{E015} \\
Qwen2.5-3B & Qwen & 3.09B & 18 & 1.200 & +0.0371 \evd{E015} \\
OPT-6.7B & OPT & 6.66B & 16 & 1.222 & +0.0233 \evd{E015} \\
Pythia-2.8B & Pythia & 2.78B & 16 & 1.240 & +0.0174 \evd{E015} \\
Qwen2.5-1.5B & Qwen & 1.54B & 14 & 1.242 & +0.0297 \evd{E015} \\
Llama-3.2-1B & Llama & 1.24B & 8 & 1.249 & +0.0360 \evd{E015} \\
OPT-2.7B & OPT & 2.65B & 16 & 1.276 & +0.0324 \evd{E015} \\
Pythia-1.4B & Pythia & 1.41B & 12 & 1.283 & +0.0213 \evd{E015} \\
Pythia-1B & Pythia & 1.01B & 8 & 1.315 & +0.0226 \evd{E015} \\
OPT-1.3B & OPT & 1.32B & 12 & 1.315 & +0.0281 \evd{E015} \\
Qwen2.5-0.5B & Qwen & 0.49B & 12 & 1.342 & +0.0356 \evd{E015} \\
GPT-2 Large & GPT-2 & 0.77B & 18 & 1.351 & +0.0173 \evd{E015} \\
Pythia-410M & Pythia & 0.41B & 12 & 1.381 & +0.0233 \evd{E015} \\
GPT-2 Medium & GPT-2 & 0.35B & 12 & 1.394 & +0.0177 \evd{E015} \\
OPT-350M & OPT & 0.33B & 12 & 1.453 & +0.0259 \evd{E015} \\
GPT-2 Small & GPT-2 & 0.12B & 6 & 1.486 & +0.0170 \evd{E015} \\
OPT-125M & OPT & 0.13B & 6 & 1.514 & +0.0147 \evd{E015} \\
Pythia-160M & Pythia & 0.16B & 6 & 1.523 & +0.0209 \evd{E015} \\
DistilGPT-2 & GPT-2 & 0.08B & 3 & 1.628 & +0.0071 \evd{E015} \\
Pythia-70M & Pythia & 0.07B & 3 & 1.699 & -0.0012 \evd{E015} \\
\bottomrule
\end{longtable}

All 22 models enter the primary middle-layer correlation.\evd{E015}
The capable-model sensitivity retains the ten models below 1.30 bits per byte.\evd{E015}
The family-factor analysis excludes Llama and Mistral because each has fewer than three model sizes; those exclusions do not apply to the primary cross-family correlation.\evd{E015}

These tables distinguish choices that changed or bounded an interpretation from checks of local stability.
Participant heterogeneity and technical-seed variation answer different questions; neither can substitute for the other or expand the claims reported in Section~\ref{sec:results}.
